% ============================================================
%  Misogyny Detection in Internet Memes
%  NLP Capstone – ACL Rolling Review format
%  Compile on Overleaf: https://www.overleaf.com
%    Template: "ACL Rolling Review" (search in Overleaf templates)
%    Replace \documentclass line if needed with the ARR version.
% ============================================================
\documentclass[11pt]{article}

% -----------------------------------------------------------
% If you upload to Overleaf with the ACL ARR template,
% replace the block below with \usepackage[review]{acl}
% -----------------------------------------------------------
\usepackage[margin=1in,columnsep=0.25in]{geometry}
\usepackage{times}
\usepackage{latexsym}
\usepackage{amsmath}
\usepackage{amssymb}
\usepackage{booktabs}
\usepackage{graphicx}
\usepackage{url}
\usepackage{microtype}
\usepackage{xcolor}
\usepackage{multirow}
\usepackage{multicol}
\usepackage[hidelinks]{hyperref}
\usepackage[T1]{fontenc}
\usepackage[utf8]{inputenc}

% Simulate ACL two-column layout
\setlength{\columnsep}{0.25in}

\title{\textbf{Misogyny Detection in Internet Memes:\\
Multi-Task Learning with Adversarial Training\\
and RAG-Based Synthetic Data Generation}}

\author{
  Wassim VQ \\
  Natural Language Processing — Capstone Project \\
  Professor: Javier Lasheras Navas \\
}

\date{}

\begin{document}
\maketitle

\begin{abstract}
We present a text-only system for detecting misogyny in internet meme captions using the MAMI dataset from SemEval-2022 Task 5 (7,500 samples). Our approach combines a domain-adapted RoBERTa backbone (pre-trained on 124M tweets) with \textbf{multi-task learning (MTL)} to jointly solve binary misogyny detection and four-label sub-type classification in a single shared model. We apply \textbf{Fast Gradient Method (FGM)} adversarial training to improve robustness to OCR noise and lexical variation, and employ per-label threshold tuning instead of a uniform 0.5 cutoff. On 1{,}500 held-out test samples, our MTL+FGM system achieves macro F1 of \textbf{0.8253} (binary) and \textbf{0.5199} (multi-label, cascaded), surpassing a TF-IDF–Logistic Regression baseline by +4.7\% and +7.7\% respectively. We complement the classifier with a complete \textbf{Retrieval-Augmented Generation (RAG)} pipeline using a locally quantized Mistral 7B model that generates type-distinguishable synthetic memes for evaluation. Error analysis reveals that 69\% of failures are attributable to implicit or visually-dependent misogyny, with negation (11\%) and irony (6\%) as the leading linguistically identifiable failure modes.
\end{abstract}

\begin{multicols}{2}

% ============================================================
\section{Introduction}
% ============================================================

Online misogyny — hostile or derogatory content targeting women — spreads rapidly through internet memes, which combine visual humor with text to carry harmful stereotypes under the guise of a joke \cite{fersini2022semeval}. Automated detection is essential for content moderation at scale, but presents distinct NLP challenges: informal register, deliberate misspellings, irony, and meaning that often depends on the image rather than the text alone.

The \textbf{MAMI} shared task at SemEval-2022 (Task 5) formalizes this problem with two sub-tasks: (A) binary detection of misogynous content, and (B) identification of sub-types — \textit{shaming}, \textit{stereotype}, \textit{objectification}, and \textit{violence}. This paper tackles both sub-tasks using a \textbf{text-only} approach on OCR-extracted captions.

Our main contributions are:
\begin{enumerate}
  \item A \textbf{multi-task model} (shared RoBERTa backbone + two classification heads) that trains both tasks jointly, giving the multi-label head access to all 7{,}500 samples through shared representations.
  \item \textbf{FGM adversarial training} applied to hate speech classification, improving robustness to the lexical noise inherent in meme text.
  \item A \textbf{cascaded prediction pipeline} enforcing the logical constraint that sub-types cannot be active unless the meme is first classified as misogynous.
  \item A full \textbf{RAG pipeline} (ChromaDB + BM25 + Mistral 7B) that generates synthetic memes used for distributional shift evaluation.
  \item Comprehensive \textbf{error analysis} with automatic failure categorization.
\end{enumerate}

% ============================================================
\section{Related Work}
% ============================================================

\paragraph{Misogyny detection.}
SemEval-2022 Task 5 attracted 83 teams; the best text-only approaches reached macro F1 $\approx$ 0.71–0.73 on Task A \cite{fersini2022semeval}. Our system surpasses this with a domain-matched pre-trained model.

\paragraph{Domain-adaptive pre-training.}
\citet{gururangan2020dont} demonstrate that continued pre-training on in-domain text consistently outperforms larger general-domain models. We apply this insight by selecting \texttt{cardiffnlp/twitter-roberta-base-hate}, pre-trained on 124M tweets and fine-tuned on HatEval and OffComEval.

\paragraph{Multi-task learning.}
\citet{ruder2017overview} surveys hard-parameter-sharing MTL, showing that auxiliary tasks act as regularizers for related tasks. \citet{caruana1997multitask} first formalized this inductive bias. Our binary task provides rich gradient signal to the backbone while the multi-label task specializes the output heads.

\paragraph{Adversarial training for NLP.}
\citet{miyato2017adversarial} proposed perturbing word embeddings with gradient-based noise, achieving improvements on text classification without any architectural change. We adopt FGM for its favorable accuracy–efficiency tradeoff over multi-step methods such as PGD \cite{madry2018towards} and FreeLB \cite{zhu2020freelb}.

\paragraph{Retrieval-Augmented Generation.}
\citet{lewis2020retrieval} introduced RAG, grounding LLM outputs in retrieved evidence. Reciprocal Rank Fusion (RRF) of dense and sparse retrievers \cite{cormack2009reciprocal} improves robustness. The ReAct framework \cite{yao2023react} enables iterative, multi-step retrieval through a reasoning–action loop.

% ============================================================
\section{Dataset and Preprocessing}
% ============================================================

\paragraph{MAMI.}
7{,}500 English meme captions annotated for: (i) binary misogyny; (ii) four sub-types (shaming, stereotype, objectification, violence). Class balance is approximately 47\%/53\% for the binary label; among misogynous memes, stereotype is the most frequent (61\%), violence the rarest (21\%).

\paragraph{Preprocessing.}
We deliberately apply \textbf{minimal preprocessing}: text columns are concatenated into a single \texttt{Text} field. No stemming, lemmatization, or stop-word removal is performed — the BPE tokenizer of RoBERTa handles sub-word variation natively, and these features are informative signal in the tweet domain.

\paragraph{Data split.}
A single stratified split (stratified by binary label, \texttt{random\_state=42}) is shared by both tasks to prevent cross-task leakage and ensure comparable metrics:

\begin{center}
\small
\begin{tabular}{lcr}
\toprule
\textbf{Partition} & \textbf{\%} & \textbf{Samples} \\
\midrule
Train & 70\% & 5{,}250 \\
Validation & 10\% & 750 \\
Test & 20\% & 1{,}500 \\
\bottomrule
\end{tabular}
\end{center}

The test set is never used for hyperparameter selection, threshold tuning, or early stopping.

% ============================================================
\section{Methodology}
% ============================================================

\subsection{Non-DL Baseline}

We implement a cascaded TF-IDF + Logistic Regression baseline following the same two-stage architecture as the transformer model to ensure fair comparison:

\begin{itemize}
  \item \textbf{Stage 1 (binary):} TF-IDF (50k features, bigrams, sublinear\_tf) + LR (C=1, class\_weight='balanced').
  \item \textbf{Stage 2 (multi-label):} \texttt{MultiOutputClassifier(LogisticRegression)} trained on misogynous-only training samples.
  \item \textbf{Cascade:} Stage 2 predictions are zeroed when Stage 1 predicts non-misogynous.
\end{itemize}

\subsection{Transformer Model Selection}

We select \texttt{cardiffnlp/twitter-roberta-base-hate} (RoBERTa-base, 125M params) over general-domain alternatives based on two domain-match criteria: (i) continued pre-training on 124M tweets matches the informal, internet-slang register of meme text; (ii) fine-tuning on HatEval and OffComEval directly transfers hate-speech representations.

Ablation experiments confirmed this choice: DistilBERT achieves binary F1 $\approx$~0.75 and DeBERTa-v3-base regresses to $\approx$~0.74 due to training instability with FGM, despite having 50\% more parameters.

\subsection{Multi-Task Learning Architecture}

\begin{center}
\small
\begin{tabular}{l}
\toprule
\textbf{MTL Model Architecture} \\
\midrule
Input $\rightarrow$ RoBERTa backbone (768-dim CLS) \\
\quad$\rightarrow$ Binary head: Drop(0.1)–Dense(768)–Tanh–Drop(0.1)–Proj(2) \\
\quad$\rightarrow$ ML head: Drop(0.1)–Dense(768)–Tanh–Drop(0.1)–Proj(4) \\
\bottomrule
\end{tabular}
\end{center}

Both heads output simultaneously from a single forward pass. The \textbf{joint loss} is:
\[
\mathcal{L} = \mathcal{L}_{\text{binary}}(\text{all}) + \mathcal{L}_{\text{ML}}(\text{misogynous only})
\]

$\mathcal{L}_{\text{binary}}$ is cross-entropy with class weights; $\mathcal{L}_{\text{ML}}$ is BCE with per-label \texttt{pos\_weight} to handle severe label imbalance (e.g., violence pos\_weight = 3.78). Using a shared backbone gives the multi-label head access to gradient signal from all 7{,}500 samples — 2.8$\times$ more than training a separate model on the misogynous-only subset (2{,}655 samples).

\subsection{FGM Adversarial Training}

At each training step, after the standard forward + backward pass, we apply one FGM perturbation step to the word embeddings:

\[
\delta = \varepsilon \cdot \frac{\nabla_e \mathcal{L}}{\|\nabla_e \mathcal{L}\|}
\]

with $\varepsilon=1.0$. A second forward pass with $e + \delta$ accumulates additional gradients. Original embeddings are then restored before the optimizer step. This simulates lexical noise (OCR errors, creative spellings) at training time without any architectural change.

\subsection{Cascaded Pipeline}

At inference, the binary head gates the multi-label head: if \texttt{binary\_pred = 0}, all sub-type predictions are forced to \textbf{[0,0,0,0]}, enforcing the logical constraint that sub-types are only meaningful for misogynous content.

\subsection{Per-Label Threshold Tuning}

After training, per-label thresholds are optimized over the misogynous-only validation subset via grid search ($n=80$ steps in [0.1, 0.9]). Tuning on misogynous-only samples avoids bias from the abundance of easy true-negatives in the full validation set.

Optimized thresholds: shaming = 0.465, stereotype = 0.100, objectification = 0.292, violence = 0.586. The wide range (0.1–0.586) confirms that a uniform 0.5 threshold is substantially suboptimal, recovering +2–5\% F1 per label.

\subsection{Optimization Details}

AdamW ($\eta=3\text{e-}5$, weight decay = 0.01) with discriminative fine-tuning (backbone $\eta/10$). Cosine schedule with 10\% linear warmup. Gradient accumulation (effective batch = 64), gradient clipping (max\_norm = 1.0), mixed-precision fp16. Early stopping (patience = 8) on combined val metric $0.4\times\text{F1}_{\text{bin}} + 0.6\times\text{F1}_{\text{ML}}$. All experiments use \texttt{RANDOM\_STATE=42} for full reproducibility.

\subsection{RAG Pipeline}

A complete Retrieval-Augmented Generation system complements the classifier:

\begin{itemize}
  \item \textbf{Ingestion:} PDF/DOCX parsing via pdfplumber/python-docx; fixed-size chunking (256 words, 32-word overlap).
  \item \textbf{Indexing:} Dense retrieval via ChromaDB (\texttt{all-MiniLM-L6-v2}) + sparse BM25.
  \item \textbf{Retrieval:} Hybrid \textbf{Reciprocal Rank Fusion (RRF)}: $\text{score}(d) = \sum_m \frac{1}{k + r_m(d)}$, $k=60$.
  \item \textbf{Generation:} Mistral 7B Instruct (Q4\_K\_M GGUF, 4.4GB) served via llama.cpp.
  \item \textbf{Bonus — ReAct agent:} Iterative SEARCH/ANSWER loop (max 5 steps) for multi-step queries.
\end{itemize}

The pipeline runs entirely locally via Docker Compose (5 services: MinIO, ChromaDB, llama.cpp, Flask API, Streamlit).

% ============================================================
\section{Experiments and Results}
% ============================================================

\subsection{Main Results}

\begin{table}[h]
\small
\centering
\caption{Test-set F1-macro (1{,}500 samples).}
\begin{tabular}{lcc}
\toprule
\textbf{System} & \textbf{Binary F1} & \textbf{ML E2E F1} \\
\midrule
TF-IDF + LR (baseline) & 0.7879 & 0.4829 \\
MTL + FGM (ours) & \textbf{0.8253} & \textbf{0.5199} \\
\midrule
$\Delta$ & +0.0374 & +0.0370 \\
Relative gain & +4.7\% & +7.7\% \\
\bottomrule
\end{tabular}
\end{table}

\begin{table}[h]
\small
\centering
\caption{Per-label results (MTL+FGM, test set, E2E cascaded).}
\begin{tabular}{lccc}
\toprule
\textbf{Label} & \textbf{P} & \textbf{R} & \textbf{F1} \\
\midrule
shaming        & 0.36 & 0.58 & 0.45 \\
stereotype     & 0.52 & 0.83 & 0.64 \\
objectification& 0.48 & 0.68 & 0.56 \\
violence       & 0.41 & 0.45 & 0.43 \\
\midrule
macro avg      & 0.44 & 0.64 & 0.52 \\
\bottomrule
\end{tabular}
\end{table}

\paragraph{Observations.}
Performance correlates with label frequency: stereotype (61\% of misogynous, F1 = 0.64) $>$ objectification (46\%, 0.56) $>$ shaming (27\%, 0.45) $>$ violence (21\%, 0.43). Recall dominates precision across all labels — a deliberate consequence of the per-label threshold optimization (e.g., stereotype threshold = 0.10), appropriate for a content moderation use case where missing harmful content is more costly than false positives.

\subsection{Ablation — Project Evolution}

\begin{table}[h]
\small
\centering
\caption{Historical model progression.}
\begin{tabular}{lcc}
\toprule
\textbf{Iteration} & \textbf{Bin F1} & \textbf{ML F1} \\
\midrule
v1: DistilBERT (separate) & 0.75 & 0.40 \\
v2: twitter-roberta (separate) & 0.8213 & 0.4999\\
\textbf{v3: MTL + FGM (shared)} & \textbf{0.8253} & \textbf{0.5199} \\
\bottomrule
\end{tabular}
\end{table}

The v1→v2 gain (+0.07 binary) is dominated by domain-matched pre-training. The v2→v3 gain is driven by MTL: the shared backbone transfers knowledge from all 7{,}500 samples to the multi-label head, producing the largest ML improvement (+0.020), while the binary gain is marginal (+0.004) since binary already had access to all samples.

\subsection{Results on RAG-Generated Memes}

\begin{table}[h]
\small
\centering
\caption{Classifier results on 50 RAG-generated memes.}
\begin{tabular}{lcc}
\toprule
& \textbf{Binary F1} & \textbf{ML E2E F1} \\
\midrule
Real test set & 0.8253 & 0.5199 \\
RAG-generated & 0.8824 & 0.4291 \\
\bottomrule
\end{tabular}
\end{table}

The higher binary F1 on generated memes (0.8824) reflects that LLM-generated captions are more \textit{explicit} than real memes, which rely on visual context and cultural implicature. The lower ML F1 (0.4291) reflects two factors: (i) generated memes are single-label by construction, diverging from the multi-label co-occurrence distribution in MAMI; (ii) small sample size (50) makes macro F1 high-variance.

\subsection{Error Analysis — Top 20 Worst Failures}

We apply automatic regex-based tagging to all misclassified test samples, identifying five failure categories.

\begin{table}[h]
\small
\centering
\caption{Failure type distribution (test set errors).}
\begin{tabular}{lrr}
\toprule
\textbf{Failure Category} & \textbf{Binary} & \textbf{ML} \\
\midrule
OTHER (implicit / visual) & 69\% (180) & 66\% (475) \\
SHORT\_TEXT ($<$8 words)   & 15\% (38)  & 14\% (101) \\
NEGATION                   & 11\% (28)  & 13\% (95)  \\
IRONY / SARCASM            & 6\% (16)   & 9\% (62)   \\
AMBIGUOUS\_HEDGE           & 1\% (3)    & 2\% (12)   \\
\midrule
\textbf{Total errors}      & \textbf{262 / 1{,}500} & \textbf{720 / 1{,}500} \\
\bottomrule
\end{tabular}
\end{table}

\paragraph{Key observations.}
\textbf{(1) OTHER dominates (67\%):} Most errors stem from \textit{implicit misogyny} — harmful content that requires cultural pragmatics, world knowledge, or the meme image to interpret. Example: \textit{"Just here to admire the scenery"} (posted with a photo of a woman) has no lexical signal yet is annotated as objectification.

\textbf{(2) SHORT\_TEXT (15\% binary):} Texts with $<$8 words provide insufficient context for the [CLS] representation. Example errors: \textit{"So true LOL"} or \textit{"Goals"} — inherently ambiguous without visual context.

\textbf{(3) NEGATION (11\% binary):} The model sometimes focuses on content words while missing negation cues. Example: \textit{"Women are definitely NOT bad at driving"} may trigger a misogynous prediction due to the co-occurrence of \textit{women}, \textit{bad}, and \textit{driving}.

\textbf{(4) IRONY/SARCASM — disproportionately high in ML (9\% vs. 6\%):} Irony is especially damaging for sub-type classification because it can simultaneously activate multiple wrong labels. Example: \textit{"Oh sure, women totally belong in the kitchen — said no functional society ever"} is sarcastic (not misogynous) but contains strong stereotype and shaming signals.

\paragraph{Limitations.}
\begin{itemize}
  \item \textbf{Text-only:} Without image features, $\approx$69\% of errors cannot be resolved by any text-based method.
  \item \textbf{Latency:} 1.06 ms/sample on GPU; 12.3 ms/sample on CPU only — suitable for batch but not real-time on CPU.
  \item \textbf{Cost:} 902.6s ($\approx$15 min) training time on GPU; no cloud API required.
  \item \textbf{Bias:} The MAMI dataset is English-only; performance on other languages or dialects is unknown. The model may inherit biases from the Twitter pre-training corpus.
\end{itemize}

% ============================================================
\section{Infrastructure and Reproducibility}
% ============================================================

The full system is containerized in a single \texttt{docker-compose.yml} (5 services). Code is structured into \texttt{/src} (model, app, backend, frontend), \texttt{/notebooks} (Labs 1–3), and \texttt{/tests} (61 unit tests, pytest). All random seeds are fixed (\texttt{RANDOM\_STATE=42}: \texttt{numpy}, \texttt{torch}, \texttt{random}, \texttt{sklearn}). Code quality: pylint 10.00/10, black formatting. The ONNX-exported model enables framework-independent inference.

% ============================================================
\section{Conclusion}
% ============================================================

We presented a production-ready NLP system for misogyny detection in internet memes combining (i) a domain-adapted RoBERTa backbone, (ii) multi-task learning for joint binary/multi-label training, (iii) FGM adversarial training, (iv) cascaded prediction enforcing logical constraints, and (v) per-label threshold tuning. The system outperforms the TF-IDF+LR baseline by +4.7\% (binary) and +7.7\% (multi-label E2E). Error analysis reveals that the dominant failure mode is implicit, visually-dependent misogyny — a fundamental limitation of text-only approaches that motivates future multimodal extensions. The complementary RAG pipeline demonstrates that the same codebase supports both discriminative classification and generative content creation at production quality.

% ============================================================
\section*{Acknowledgments}
% ============================================================

We thank Professor Javier Lasheras Navas for the project guidelines and the Cardiff NLP group for releasing \texttt{twitter-roberta-base-hate}.

% ============================================================
\begin{thebibliography}{10}

\bibitem{fersini2022semeval}
Elisabetta Fersini, Francesca Gasparetto, and Aurora Savini.
\newblock {SemEval-2022 Task 5: Multimedia Automatic Misogyny Identification}.
\newblock In \textit{Proceedings of the 16th International Workshop on Semantic Evaluation (SemEval-2022)}, 2022.

\bibitem{gururangan2020dont}
Suchin Gururangan, Ana Marasović, Swabha Swayamanta, Kyle Lo, Iz Beltagy, Doug Downey, and Noah A. Smith.
\newblock Don't stop pretraining: Adapt language models to domains and tasks.
\newblock In \textit{Proceedings of ACL}, 2020.

\bibitem{ruder2017overview}
Sebastian Ruder.
\newblock An overview of multi-task learning in deep neural networks.
\newblock \textit{arXiv preprint arXiv:1706.05098}, 2017.

\bibitem{caruana1997multitask}
Rich Caruana.
\newblock Multitask learning.
\newblock \textit{Machine Learning}, 28(1), 1997.

\bibitem{miyato2017adversarial}
Takeru Miyato, Andrew M. Dai, and Ian Goodfellow.
\newblock Adversarial training methods for semi-supervised text classification.
\newblock In \textit{Proceedings of ICLR}, 2017.

\bibitem{madry2018towards}
Aleksander Madry, Aleksandar Makelov, Ludwig Schmidt, Dimitris Tsipras, and Adrian Vladu.
\newblock Towards deep learning models resistant to adversarial attacks.
\newblock In \textit{Proceedings of ICLR}, 2018.

\bibitem{zhu2020freelb}
Chen Zhu, Yu Cheng, Zhe Gan, Siqi Sun, Tom Goldstein, and Jingjing Liu.
\newblock {FreeLB}: Enhanced adversarial training for natural language understanding.
\newblock In \textit{Proceedings of ICLR}, 2020.

\bibitem{lewis2020retrieval}
Patrick Lewis, Ethan Perez, Aleksandra Piktus, et al.
\newblock Retrieval-augmented generation for knowledge-intensive {NLP} tasks.
\newblock In \textit{Proceedings of NeurIPS}, 2020.

\bibitem{cormack2009reciprocal}
Gordon V. Cormack, Charles L.A. Clarke, and Stefan Buettcher.
\newblock Reciprocal rank fusion outperforms condorcet and individual rank learning methods.
\newblock In \textit{Proceedings of SIGIR}, 2009.

\bibitem{yao2023react}
Shunyu Yao, Jeffrey Zhao, Dian Yu, Nan Du, Izhak Shafran, Karthik Narasimhan, and Yuan Cao.
\newblock {ReAct}: Synergizing reasoning and acting in language models.
\newblock In \textit{Proceedings of ICLR}, 2023.

\end{thebibliography}
% ============================================================

\end{multicols}
\end{document}
